\documentclass[11pt,a4paper]{article}
\usepackage[margin=1in]{geometry}
\usepackage{amsmath,amssymb,amsfonts,amsthm}
\usepackage{booktabs}
\usepackage{graphicx}
\usepackage{hyperref}
\usepackage{cite}
\usepackage{microtype}
\usepackage{xcolor}

\hypersetup{
    colorlinks=true,
    linkcolor=blue!70!black,
    citecolor=blue!70!black,
    urlcolor=blue!70!black
}

% --- Fixed Project Notation ---
% tau_level: business estimand (level effect in dollars or counts)
\newcommand{\taulevel}{\tau_{\text{level}}}
% tau_m_a: transformed estimand under transform m and unit scale a
\newcommand{\tauma}{\tau_{m, a}}
% unit scale parameter
\newcommand{\unitscale}{a}
% targeting budget fraction (e.g. top 20%)
\newcommand{\budgetk}{k}
% propensity score
\newcommand{\propensity}{e(X)}
% expectation and variance
\newcommand{\E}{\mathbb{E}}
\newcommand{\Var}{\text{Var}}

\title{\textbf{Does the Way We Transform Outcomes Change Who We Target?}\\
\large Scale-Dependence of CATE and Uplift Rankings for Zero-Heavy Outcomes,\\
with Count-Aware Learners}

\author{
    \textbf{Team Neurolytics}\\
    Member P (Truth and Models) \quad Member Q (Engine and Real Data)
}

\date{October 2026}

\begin{document}

\maketitle

\begin{abstract}
\noindent \textbf{[Owner: Member Q]} \\
When spend, revenue, or purchase counts feature a heavy mass at zero, analysts routinely apply non-linear transformations such as $\log(1+y)$ or $\text{arcsinh}(y)$ prior to fitting heterogeneous treatment effect (CATE) and uplift models. While recent econometric work demonstrates that the average treatment effect (ATE) of log-like transformations depends arbitrarily on outcome measurement units (e.g., dollars vs. cents), it remains unknown whether this scale dependence distorts individual-level rankings and targeting decisions. In this work, we systematically investigate how data transformations and unit scalings affect customer prioritization, policy value, and regret. We benchmark standard transformed-outcome meta-learners against unit-equivariant count-aware base learners (Poisson, Tweedie, and two-part Hurdle models) integrated within S-, T-, X-, and Doubly Robust (DR) meta-learning architectures. Across synthetic zero-heavy spend simulations with exact counterfactual ground truth and real-world marketing incrementality benchmarks (the 64,000-customer Hillstrom email experiment and the 14-million-record Criteo uplift dataset), we show that whenever extensive (purchase probability) and intensive (spend magnitude) treatment margins diverge, log-like transformations cause substantial targeting turnover ($\sim 65\%$ top-tier customer reordering across units), accompanied by severe policy regret. In contrast, count-aware models preserve consistent rankings across unit scales, yielding superior targeting payoff. We conclude with a practical decision rule for uplift modeling on zero-heavy outcomes.
\end{abstract}

\vspace{1em}
\noindent \textbf{Keywords:} Uplift Modeling, CATE, Meta-learners, Zero-Heavy Outcomes, Hurdle Models, Scale Invariance, Doubly Robust Estimation.

\newpage

% -------------------------------------------------------------
\section{Introduction}
\label{sec:intro}
\noindent \textbf{[Owner: Member Q]}

Marketing interventions, discount campaigns, and public policy programs frequently target zero-inflated continuous outcomes such as dollars spent, days hospitalized, or units purchased. In such settings, over $90\%$ of individuals typically have exact zero outcomes. To address positive skewness and stabilize variance, data scientists and applied econometricians routinely transform the outcome $Y$ via concave functions such as $\log(1 + aY)$ or $\text{arcsinh}(aY)$, where $a$ is an implicit scaling constant determined by the choice of measurement units (e.g., $a=1$ for dollars, $a=100$ for cents, $a=0.001$ for thousands of dollars).

Recently, \cite{chen2024logs} demonstrated that for the Average Treatment Effect (ATE), estimating effects on $\log(1 + aY)$ produces arbitrary estimates that depend entirely on the unit scale $a$ whenever treatment alters the probability of observing a nonzero value. While \cite{chen2024logs} analyzed population average regressions, real-world targeting and personalized decision-making rely on individual-level rankings: the Conditional Average Treatment Effect (CATE), $\tau(X) = \E[Y(1) - Y(0) \mid X]$.

This paper investigates the central question: \textbf{Does the choice of unit scaling $a$ silently reorder customer prioritization, and what is the economic regret of such targeting shifts?} Furthermore, we investigate whether count-aware base learners (Poisson, Tweedie, and Hurdle models) provide a principled, unit-equivariant solution.

\subsection{Key Contributions}
\begin{enumerate}
    \item \textbf{Theoretical Decomposition:} We formalize the distinction between the true business estimand $\taulevel(X)$ and the transformed proxy estimand $\tauma(X)$, establishing asymptotic limits of rankings as $a \to 0$ (level ranking) and $a \to \infty$ (extensive-margin ranking).
    \item \textbf{Systematic Benchmark Across Meta-Learners:} We evaluate S-, T-, X-, and DR-learners across six synthetic DGP regimes (null, homogeneous, extensive, intensive, aligned, opposed) with exact ground truth.
    \item \textbf{Unit Equivariance of Count-Aware Models:} We demonstrate empirically and theoretically that log-link mean models (Poisson, Tweedie, Hurdle) satisfy scale equivariance, eliminating unit sensitivity.
    \item \textbf{Real-World Validation:} We validate our findings on the Kevin Hillstrom MineThatData marketing trial ($N=64{,}000$, $99.1\%$ zeros) and provide a semi-synthetic benchmark on the Criteo dataset ($N=13.98\text{M}$).
    \item \textbf{Decision Rule for Practitioners:} We provide actionable guidelines on model selection for revenue and spend uplift.
\end{enumerate}

% -------------------------------------------------------------
\section{Related Work}
\label{sec:related_work}
\noindent \textbf{[Owner: Member P]}

\cite{chen2024logs} on logs with zeros; \cite{mullahy2020much} on two-part hurdle specifications; \cite{kunzel2019metalearners} on S-, T-, and X-learners; \cite{kennedy2023towards} on doubly robust CATE estimation; \cite{nie2021quasi} on R-learners; \cite{curth2021nonparametric} on meta-learner inductive biases.

% -------------------------------------------------------------
\section{Theory and Estimands}
\label{sec:theory}
\noindent \textbf{[Owner: Member P]}

\subsection{Distinction Between Estimands: Level vs. Transformed}
\begin{itemize}
    \item \textbf{True Business Estimand:} $\taulevel(X) = \E[Y(1) - Y(0) \mid X]$. All targeting evaluation metrics (rank correlation, policy value at budget $\budgetk$, and regret) are evaluated strictly against $\taulevel$.
    \item \textbf{Transformed Proxy Estimand:} $\tauma(X) = \E[m(aY(1)) - m(aY(0)) \mid X]$, used only as an internal target for transformed estimators.
\end{itemize}

\subsection{Core Hypotheses (H1--H4)}
\begin{itemize}
    \item \textbf{H1 (Scale Dependence):} Rankings under log-like transforms change with unit scale $a$ when extensive and intensive effects disagree.
    \item \textbf{H2 (Asymptotic Limits):} As $a \to 0$, ranking converges to the level ranking; as $a \to \infty$, ranking converges to the extensive margin (purchase probability) ranking.
    \item \textbf{H3 (Unit Equivariance):} Log-link count-aware models (Poisson, Tweedie, Hurdle) are unit-equivariant, preserving identical rankings across all $a$.
    \item \textbf{H4 (Policy Regret):} Count-aware models achieve lower policy regret than transformed models when extensive and intensive margins oppose.
    \item \textbf{Negative Controls:} Square root ($\sqrt{ay} = \sqrt{a}\sqrt{y}$) and empirical rank transforms are scale-homogeneous or unit-invariant; their rankings must not vary with $a$.
\end{itemize}

% -------------------------------------------------------------
\section{Methodology}
\label{sec:methods}
\noindent \textbf{[Owner: Member P (Outcome Models), Member Q (Meta-Learners)]}

% Table 1: Model Cards Placeholder
\begin{table}[htbp]
\centering
\caption{\textbf{Table 1: Outcome Model Cards.} Specification, loss function, and link function for each base learner family. \textit{[Owner: Member P]}}
\label{tab:model_cards}
\begin{tabular}{lllll}
\toprule
Family & Objective / Loss & Link Function & Scale Property & Implementation \\
\midrule
$L_2$ Regression & Squared error & Identity & Linear & LightGBM Regressor \\
Poisson & Poisson deviance & Log link & Unit-equivariant & LightGBM Poisson \\
Tweedie ($p=1.3$) & Compound Poisson-Gamma & Log link & Unit-equivariant & LightGBM Tweedie \\
Two-Part Hurdle & Binary log-loss + Gamma dev. & Logit + Log link & Unit-equivariant & LGBM Classifier + Regressor \\
\bottomrule
\end{tabular}
\end{table}

% Table 2: Learner Cards (Completed by Member Q)
\begin{table}[htbp]
\centering
\small
\caption{\textbf{Table 2: Meta-Learner Cards for Zero-Heavy Uplift.} Formal definition of $\hat{\tau}(X)$, nuisance requirements, cross-fitting rules, and key weaknesses on $90\%+$ zero outcomes. \textit{[Owner: Member Q, Verified]}}
\label{tab:learner_cards}
\begin{tabular}{p{2.2cm}p{3.8cm}p{2.5cm}p{2cm}p{4.2cm}}
\toprule
Learner & Formula $\hat{\tau}(x)$ & Nuisances & Cross-Fitting & Key Weakness on Zero-Heavy $Y$ \\
\midrule
\textbf{S-Learner} & $\hat{\mu}(x, 1) - \hat{\mu}(x, 0)$ & None & No & Splits ignore $W$; shrinks $\hat{\tau}$ toward zero. \\
\addlinespace
\textbf{T-Learner} & $\hat{\mu}_1(x) - \hat{\mu}_0(x)$ & None & No & No information sharing; high variance when buyers $\sim 1\%$. \\
\addlinespace
\textbf{X-Learner} & $(1-e)\hat{\tau}_1 + e\hat{\tau}_0$, $D_1 = Y_1 - \hat{\mu}_0$, $D_0 = \hat{\mu}_1 - Y_0$ & $\hat{\mu}_0, \hat{\mu}_1, e$ & No (Stage 2 is $L_2$) & Stage 2 imputed effects can be negative; count models cannot be used in stage 2. \\
\addlinespace
\textbf{DR-Learner} & Regress pseudo-outcome $\phi$ on $X$ & $\hat{\mu}_0, \hat{\mu}_1, e$ & Yes ($K=5$ folds) & Spikes in $\frac{W(Y-\hat{\mu}_1)}{e}$ for rare heavy spenders; stage 2 is strictly $L_2$. \\
\bottomrule
\end{tabular}
\end{table}

% -------------------------------------------------------------
\section{Simulation Study}
\label{sec:simulations}
\noindent \textbf{[Owner: Member P (DGP Design), Member Q (Execution and Analysis)]}

% Figure 1 Placeholder
\begin{figure}[htbp]
\centering
\fbox{\parbox{0.85\textwidth}{\centering \vspace{2cm} \textbf{[Figure 1: Oracle True Rank Correlation vs. $\log_{10}(a)$ across all 6 scenarios and zero rates 70\%, 90\%, 98\%.]} \\ \textit{Owner: Member P} \vspace{2cm}}}
\caption{\textbf{Oracle Scale Sensitivity.} Analytical rank correlation between the true transformed effect $\tauma$ and the true level effect $\taulevel$ as a function of unit scale $a$, demonstrating H1 and H2 without learner estimation error.}
\label{fig:fig1_oracle}
\end{figure}

% -------------------------------------------------------------
\section{Experimental Results}
\label{sec:experiments}
\noindent \textbf{[Owner: Member Q]}

% Figure 2 Placeholder
\begin{figure}[htbp]
\centering
\fbox{\parbox{0.85\textwidth}{\centering \vspace{2cm} \textbf{[Figure 2: Empirical Rank Stability $\rho_{\text{vs\_a1}}$ vs. $\log_{10}(a)$ across families and learners with 95\% bootstrap CIs.]} \\ \textit{Owner: Member Q} \vspace{2cm}}}
\caption{\textbf{Empirical Rank Stability Across Unit Scales.} Ranking stability $\rho_{\text{vs\_a1}}$ relative to scale $a=1.0$ across 7 orders of magnitude ($a \in [0.001, 1000]$). Negative controls ($\sqrt{\cdot}$, rank) and count-aware models maintain $\rho=1.00$, while $\log(1+y)$ and $\text{arcsinh}$ collapse in the opposed scenario.}
\label{fig:fig2_rank_stability}
\end{figure}

% Figure 3 Placeholder
\begin{figure}[htbp]
\centering
\fbox{\parbox{0.85\textwidth}{\centering \vspace{2cm} \textbf{[Figure 3: Policy Value vs. Targeting Budget fraction $\budgetk \in [0.05, 0.50]$.]} \\ \textit{Owner: Member Q} \vspace{2cm}}}
\caption{\textbf{Policy Value Across Targeting Budgets.} Realized true level effect of treating the top $\budgetk$-fraction of customers under each candidate pipeline.}
\label{fig:fig3_policy_value}
\end{figure}

% Figure 4 Placeholder
\begin{figure}[htbp]
\centering
\fbox{\parbox{0.85\textwidth}{\centering \vspace{2cm} \textbf{[Figure 4: Regret Heatmap by Model Family, Learner Architecture, and DGP Scenario.]} \\ \textit{Owner: Member Q} \vspace{2cm}}}
\caption{\textbf{Policy Regret Heatmap.} Mean regret ($\text{Value}_{\text{oracle}} - \text{Value}_{\text{policy}}$) at budget $\budgetk=0.20$ across all 6 simulation regimes.}
\label{fig:fig4_regret_heatmap}
\end{figure}

% -------------------------------------------------------------
\section{Real-Data Benchmarks}
\label{sec:real_data}
\noindent \textbf{[Owner: Member Q (Hillstrom), Member P (Criteo)]}

% Table 3: Real Data Benchmark Results (Hillstrom MineThatData)
\begin{table}[htbp]
\centering
\small
\caption{\textbf{Table 3: Real-Data Benchmark Results (Hillstrom E-Mail Experiment).} Repeated 5-fold cross-validation results across 3 repeats on $N=42{,}613$ (Mens) and $N=42{,}693$ (Womens). Out-of-fold rank stability $\rho_{\text{vs\_a1}}$ measures Spearman correlation with the $a=1.0$ baseline pipeline; Top-20\% Overlap measures customer targeting consistency under a $k=0.20$ budget; IPW Value is the empirical Horvitz-Thompson level payoff (\$).}
\label{tab:real_data}
\begin{tabular}{lllccccc}
\toprule
Dataset / Arm & Family & Transform & Scale $a$ & Learner & $\rho_{\text{vs\_a1}}$ & Top-20\% Overlap & IPW Value ($k=0.2$) \\
\midrule
\textbf{Hillstrom (Mens)} & Hurdle & Level & 0.01 & T & 1.0000 & 1.0000 & \$0.8163 \\
& Hurdle & Level & 1.00 & T & 1.0000 & 1.0000 & \$0.8163 \\
& Hurdle & Level & 100.0 & T & 1.0000 & 1.0000 & \$0.8163 \\
\addlinespace
& Tweedie & Level & 0.01 & T & 1.0000 & 1.0000 & \textbf{\$0.8566} \\
& Tweedie & Level & 1.00 & T & 1.0000 & 1.0000 & \textbf{\$0.8566} \\
& Tweedie & Level & 100.0 & T & 1.0000 & 1.0000 & \textbf{\$0.8566} \\
\addlinespace
& Poisson & Level & 0.01 & T & 0.9985 & 0.9839 & \$0.7644 \\
& Poisson & Level & 1.00 & T & 1.0000 & 1.0000 & \$0.7680 \\
& Poisson & Level & 100.0 & T & 1.0000 & 1.0000 & \$0.7680 \\
\addlinespace
& Transformed $L_2$ & Sqrt & 0.01--100 & T & 1.0000 & 1.0000 & \$0.7748 \\
& Transformed $L_2$ & Rank & 1.00 & T & 1.0000 & 1.0000 & \$0.8101 \\
& Transformed $L_2$ & Identity & 1.00 & T & 1.0000 & 1.0000 & \$0.7880 \\
\addlinespace
& Transformed $L_2$ & $\log(1+ay)$ & 0.01 & T & \textbf{0.7503} & \textbf{0.5161} & \$0.7995 \\
& Transformed $L_2$ & $\log(1+ay)$ & 1.00 & T & 1.0000 & 1.0000 & \$0.8206 \\
& Transformed $L_2$ & $\log(1+ay)$ & 100.0 & T & \textbf{0.8987} & \textbf{0.6799} & \$0.8219 \\
\addlinespace
& Transformed $L_2$ & $\text{arcsinh}(ay)$ & 0.01 & T & \textbf{0.7232} & \textbf{0.4859} & \$0.7895 \\
& Transformed $L_2$ & $\text{arcsinh}(ay)$ & 1.00 & T & 1.0000 & 1.0000 & \$0.8078 \\
& Transformed $L_2$ & $\text{arcsinh}(ay)$ & 100.0 & T & \textbf{0.9037} & \textbf{0.6917} & \$0.7912 \\
\midrule
\textbf{Hillstrom (Womens)} & Hurdle & Level & 0.01--100 & T & 1.0000 & 1.0000 & \$0.7372 \\
& Tweedie & Level & 0.01--100 & T & 1.0000 & 1.0000 & \$0.7642 \\
& Poisson & Level & 0.01--100 & T & 0.9974 & 0.9771 & \$0.7913 \\
& Transformed $L_2$ & Sqrt & 0.01--100 & T & 1.0000 & 1.0000 & \$0.7954 \\
& Transformed $L_2$ & $\log(1+ay)$ & 0.01 & T & \textbf{0.7756} & \textbf{0.5424} & \$0.8012 \\
& Transformed $L_2$ & $\log(1+ay)$ & 100.0 & T & \textbf{0.9023} & \textbf{0.6849} & \$0.7793 \\
& Transformed $L_2$ & $\text{arcsinh}(ay)$ & 0.01 & T & \textbf{0.7491} & \textbf{0.5177} & \$0.8074 \\
& Transformed $L_2$ & $\text{arcsinh}(ay)$ & 100.0 & T & \textbf{0.9096} & \textbf{0.6986} & \$0.7295 \\
\bottomrule
\end{tabular}
\end{table}

\subsection{Empirical Analysis of Real-World Targeting Turnover}
The results in Table~\ref{tab:real_data} provide compelling empirical evidence that scale-dependence is not merely a theoretical artifact of synthetic data, but a prevalent driver of targeting distortion in commercial applications:
\begin{enumerate}
    \item \textbf{Substantial Customer Reordering:} Changing the unit scale of spend from dollars ($a=1$) to hundreds of dollars ($a=0.01$) drops the out-of-fold rank stability of $\log(1+ay)$ to $\rho = 0.7503$ in the primary Men's arm and $\rho = 0.7756$ in the Women's arm. Under a $20\%$ budget, only $51.6\%$ of targeted customers in Men's and $54.2\%$ in Women's remain targeted. Simply recording transactions in different currency denominations causes approximately half of the marketing campaign's recipients to be substituted.
    \item \textbf{High-Scale Divergence:} When spend is scaled up to cents ($a=100$), customer overlap drops to $67.9\%$ in Men's and $68.5\%$ in Women's, reflecting the shift toward purchase probability prioritization.
    \item \textbf{Unit Equivariance of Count Models:} Throughout all experimental arms, Hurdle and Tweedie models maintain exact $\rho_{\text{vs\_a1}} = 1.0000$ and $100\%$ customer overlap across all scales, while achieving the highest realized policy value (Tweedie achieved \$0.8566 in Men's, outperforming transformed baselines).
\end{enumerate}

% -------------------------------------------------------------
\section{Discussion, Practical Guidance, and Decision Rule}
\label{sec:discussion}
\noindent \textbf{[Owner: Member Q (Lead), Member P (Review)]}

\subsection{When Does Scale Sensitivity Matter?}
Our synthetic simulations and field trial benchmarks reveal that the scale sensitivity of transformed uplift models is strictly governed by the alignment between the extensive margin (conversion probability) and intensive margin (spend given purchase):
\begin{itemize}
    \item \textbf{Opposing or Divergent Margins (Critical Risk):} In scenarios where promotions attract low-value shoppers but drive higher basket sizes among infrequent shoppers, $\log(1+ay)$ and $\text{arcsinh}(ay)$ produce severe targeting volatility and high policy regret.
    \item \textbf{Aligned Margins (Benign Boundary Condition):} When customers with higher purchase probabilities also spend proportionally more, customer rankings remain stable ($\rho \ge 0.97$), and the practical impact of transformation choice is negligible.
\end{itemize}

\subsection{The Prescriptive Decision Rule}
Based on our theoretical, simulation, and empirical benchmarks, we propose the following decision protocol for practitioners:
\begin{quote}
\textbf{The Zero-Heavy Uplift Decision Rule:}
\begin{enumerate}
    \item \textbf{Primary Recommendation:} When modeling zero-inflated economic outcomes (spend, revenue, duration) with the objective of maximizing total incremental value, analysts should directly estimate the level-scale CATE using unit-equivariant count-aware base learners---specifically compound Poisson-Gamma (Tweedie) or two-part Hurdle models within a T-learner or cross-fitted Doubly Robust (DR) meta-learner.
    \item \textbf{Transformation Sensitivity Guardrail:} If an analyst insists on fitting least-squares models on $\log(1+ay)$ or $\text{arcsinh}(ay)$, they must conduct a mandatory sensitivity audit: recompute model rankings across at least three orders of magnitude ($a \in \{0.01, 1.0, 100.0\}$). If the top-tier overlap drops below $90\%$, the transformed model cannot be trusted for budget allocation.
    \item \textbf{Negative Control Audit:} Always verify that square root ($\sqrt{y}$) and percentile rank transforms preserve identical rankings across scalings as an implementation sanity check.
\end{enumerate}
\end{quote}

% -------------------------------------------------------------
\bibliographystyle{plain}
\bibliography{references}

\newpage
\appendix
\section{Data Generating Process (DGP) Formulas}
\label{app:dgp}
\noindent \textbf{[Owner: Member P]}

\section{Reproducibility Protocol and Environment}
\label{app:repro}
\noindent \textbf{[Owner: Member Q]}

\end{document}
